% Werke X/1 (1917), native142–144. Historical note by Felix Klein.
\clearpage
\SourceLeaf{142}\NoSourceHead
\pdfbookmark[0]{Analysis}{analysis}
\vspace*{2in}
\HistoricTitle{ANALYSIS.}
\begin{center}NACHTRÄGE ZU DEN BÄNDEN III UND VIII.\end{center}
\clearpage
\SourceLeaf{143}
\SourceLeaf{144}\ArithmeticHead{137}{R}{}
\begin{center}\bfseries
NACHTRÄGLICHE BEMERKUNG ZU DER S. 36--64 DES VIII. BANDES ABGEDRUCKTEN
ABHANDLUNG: DE INTEGRATIONE FORMULAE DIFFERENTIALIS $(1+n\cos\varphi)^\nu d\varphi$.
\end{center}
\noindent\hspace*{\parindent}Es ist hier ein bedauernswerter Irrtum zu berichtigen, der sich bei der früheren Herausgabe von
Stücken des \textsc{Gauss}schen Nachlasses, die sich auf Analysis beziehen, eingeschlichen hat (wie ich dies
übrigens bereits 1902 in meinem damaligen \textit{Bericht über den Stand der Herausgabe von Gauss’ Werken},
Nachrichten von der K. Gesellschaft der Wissenschaften zu Göttingen, Geschäftliche Mitteilungen S. 12,
ausgeführt habe).

Im Jahre 1893 wurde von \textsc{Wilhelm Meyer} (Göttingen) in den Akten der Königlichen Gesellschaft
der Wissenschaften ein Manuskript mit dem Titel \textit{De integratione formulae differentialis} $(1+n\cos\varphi)^\nu d\varphi$
gefunden, das weder einen Verfassernamen noch eine Angabe über Ort und Zeit der Abfassung enthielt,
dessen Handschrift aber derjenigen des jugendlichen \textsc{Gauss} so ähnlich war, daß über den Verfasser kein
Zweifel zu bestehen schien. So wurde es von E. \textsc{Schering} in den Nachrichten von der K. Gesellschaft der
Wissenschaften vom Jahre 1893, S. 617--646, als von \textsc{Gauss} herrührend veröffentlicht und auch im Bd. VIII
der Werke, 1900, S. 35--64 wieder abgedruckt, trotzdem R. \textsc{Fricke} (der im Bd. VIII die Beiträge zur Analysis bearbeitete)
auf eine Reihe von Unstimmigkeiten aufmerksam wurde, denen er in der Anmerkung auf
S. 64 des genannten Bandes Ausdruck gegeben hat.

Ganz unabhängig davon hat mir bald darauf \textsc{Conrad Müller}, der sich damals auf meine Veranlassung
mit Studien über die an der Göttinger Universität um die Wende des 18. Jahrhunderts wirkenden
Mathematiker beschäftigte, eines Tages (1902) berichtet, daß er in den Göttingischen Anzeigen von 1800,
Stück 31, eine Besprechung von \textsc{Kästner} über eine von \textsc{Thibaut} eingereichte Schrift gefunden habe, die
im wesentlichen eine vom Verfasser selbst herrührende Inhaltsangabe enthalte. Die Vergleichung ergab
dann sofort, daß es sich dabei genau um die vorgenannte, also fälschlich \textsc{Gauss} zugeschriebene Arbeit
handelte. Nun ist wirklich jeder Zweifel ausgeschlossen: nicht \textsc{Gauss}, sondern \textsc{Thibaut} ist der Verfasser.
\par\noindent\hfill\textsc{Klein}.\par
\ClosingRule
\vfill\noindent\small X1.\hfill18\hspace*{.15\linewidth}\normalsize
\clearpage
